\documentclass[10pt,a4paper]{amsart}
\usepackage[margin=19mm]{geometry}
\usepackage{amsmath,amssymb,mathtools}
\usepackage[scr=rsfs,cal=euler]{mathalfa}
\usepackage{xfrac}
\usepackage[unicode,hidelinks,bookmarks=false]{hyperref}
\newcommand{\ie}{i.e.\ }
\newcommand{\cf}{cf.\ }
\newcommand{\resp}{respectively\ }
\newcommand{\Subsection}{\subsection}
\providecommand{\nobreakdash}{\nobreak\hbox{-}}
\setlength{\emergencystretch}{2em}
\multlinegap0pt
\usepackage{bm}
\usepackage{bbm}
\usepackage{dsfont}
\usepackage{xspace}
\usepackage{cancel}
\usepackage{mathtools}
\usepackage{amsthm}
\makeatletter
\@ifundefined{definition}{\newtheorem{definition}{Definition}}{}
\@ifundefined{lemma}{\newtheorem{lemma}[definition]{Lemma}}{}
\@ifundefined{theorem}{\newtheorem{theorem}[definition]{Theorem}}{}
\makeatother
\newcommand{\Z}{\mathbb{Z}} 
\newcommand{\cc}{\overline} 
\newcommand{\del}{\partial} 
\newcommand{\tn}{\textnormal} 
\newcommand{\vareps}{\varepsilon} 
\renewcommand{\vec}[1]{\hat{\boldsymbol{#1}}} 
\title[Steklov eigenvalues of nearly hyperspherical domains]{Steklov eigenvalues of nearly hyperspherical domains}
\author{Chee Han Tan, Robert Viator}
\date{Source: arXiv:2310.03960v1 (CC BY 4.0)}
\begin{document}
\maketitle

\noindent\emph{Source and licence.} Steklov eigenvalues of nearly hyperspherical domains. Version arXiv:2310.03960v1 (CC BY 4.0), distributed under the Creative Commons Attribution 4.0 International licence, as recorded in the arXiv licence metadata for this exact version and shown on the abstract page https://arxiv.org/abs/2310.03960v1. Full rights notice: licenses/arxiv-2310.03960-CC-BY-4.0.txt.\par\smallskip

\noindent\textit{Editorial scope.} Counted unit: the Lemma whose source label is \texttt{thm:Matrix\_triple} in the article (source0.tex lines 558--567, source label \texttt{thm:Matrix\_triple}), delivered together with the article's own complete original proof of it (source0.tex lines 568--598, one \texttt{proof} environment of 31 lines). No mathematics is written, repaired or re-derived: statement and proof are the article's own text line for line. Required context retained uncounted: eq:Domain (lines 114--116); eq:Steklov\_multiplicity (lines 118--120); eq:SCgrad (lines 185--187); eq:SCbeltrami (lines 193--195); eq:SH\_eigs (lines 213--215); eq:Oeps\_matrix (lines 398--403); thm:Oeps\_matrix (lines 411--413); eq:Oeps\_matrix1 (lines 433--435); eq:Oeps\_matrix1b (lines 454--461). Every definition, display and label of those lines is retained unchanged. Omitted from this package and not redistributed: 1--113, 117--117, 121--184, 188--192, 196--212, 216--397, 404--410, 414--432, 436--453, 462--557, 599--792. Nothing inside a retained excerpt is omitted or altered: every retained body is delivered exactly as the article's lines are written, so no omission receipt of any kind accompanies this package. Citations: the retained excerpts carry no citation command at all, so no bibliography entry of the article is transcribed and no reference is redistributed. Reference adaptation: none was needed and none was made; every reference inside the delivered unit resolves inside it, so no unresolved reference or citation is produced. Printed slips kept literal: no printed word, symbol, hypothesis or number of the counted statement or of its proof was altered. Layout and class: the article's own class is \texttt{amsart} with packages including fullpage, amsaddr, graphicx, amsmath, amssymb, amsfonts, amsthm, mathtools, enumerate, color, url, hyperref, biblatex; the wrapper is the lane's pinned \texttt{amsart} editorial wrapper at 10pt on A4, which restates the theorem-like environments the retained text needs and adds the article's own macro definitions that the retained text needs (\texttt{\textbackslash Z}, \texttt{\textbackslash cc}, \texttt{\textbackslash del}, \texttt{\textbackslash tn}, \texttt{\textbackslash vareps}, \texttt{\textbackslash vec}), with extra wrapper packages (bm, bbm, dsfont, xspace, cancel, mathtools). The delivered numbering is the unit's own. Assets: no retained span contains a figure environment, a graphics-inclusion command, a PDF-page inclusion, a table or a TikZ picture; no image file is copied into this package, the package declares no asset dependency and no image file is redistributed with it. Licence: version arXiv:2310.03960v1 of this article is distributed under the Creative Commons Attribution 4.0 International licence, as recorded in the arXiv licence metadata for this exact version and shown on the abstract page https://arxiv.org/abs/2310.03960v1. A full rights notice is delivered at licenses/arxiv-2310.03960-CC-BY-4.0.txt. Characters: every retained excerpt is pure ASCII, so the delivered engine, XeLaTeX with its default Latin Modern text font, reports no missing character.
\begin{equation} \label{eq:Domain} 
\Omega_\vareps = \left\{(r, \hat\theta): 0\le r\le 1+ \vareps\rho(\hat\theta), \, \hat\theta \in S^d \right\}, \ \ \rho(\hat\theta) = \sum_{p = 0}^\infty\sum_{q = 1}^{N(d, p)} A_{p, q} Y_{p, q}(\hat\theta). 
\end{equation} 

\begin{equation} \label{eq:Steklov_multiplicity} 
N(d, 0) = 1 \ \ \tn{ and } \ \ N(d, \ell) = \binom{d + \ell}{d} - \binom{d + \ell - 2}{d}, \ \ \ell\ge 1. 
\end{equation} 

\begin{equation} \label{eq:SCgrad} 
\nabla_{S^d} = \sum_{j = 1}^d \frac{1}{\eta_j}\frac{\del}{\del\theta_j}\vec{\theta}_j. 
\end{equation} 

\begin{equation} \label{eq:SCbeltrami} 
\Delta_{S^d} = \sum_{j = 1}^d \frac{1}{\eta_j^2\sin^{j - 1}(\theta_j)}\frac{\del}{\del\theta_j}\left(\sin^{j - 1}(\theta_j)\frac{\del}{\del\theta_j}\right). 
\end{equation} 

\begin{equation} \label{eq:SH_eigs} 
\Delta_{S^d}Y_\ell^m(\hat\theta) = -\ell(\ell + d - 1)Y_\ell^m(\hat\theta), \ \ 1\le m\le N(d, \ell). 
\end{equation} 

\begin{gather} \label{eq:Oeps_matrix}
\begin{aligned}  
M_{m, n}^{(d, k)} & = -\int_{S^d} k\rho Y_k^m\cc{Y_k^n}\, d\sigma_d - \int_{S^d} \left(\nabla_{S^d}\rho\cdot \nabla_{S^d}Y_k^m\right)\cc{Y_k^n}\, d\sigma_d \\ 
% & \stackrel{\eqref{eq:SCgrad}}{=} -\int_{S^d} k\rho Y_k^m\cc{Y_k^n}\, d\sigma_d - \sum_{j = 1}^d \int_{S^d} \frac{\del_j\rho}{\eta_j^2}\del_jY_k^m\cdot \cc{Y_k^n}\, d\sigma_d.
\end{aligned} 
\end{gather} 

\begin{theorem} \label{thm:Oeps_matrix} 
Given $d\ge 3$ and $k\in\Z^+$, define $N_k = \sum_{\ell = 1}^{k - 1} N(d, \ell)$, where $N(d, \ell)$ is defined by \eqref{eq:Steklov_multiplicity}. For $N_k + 1\le n\le N_{k + 1}$, the Steklov eigenvalues $\lambda_n(\vareps)$ of a nearly hyperspherical domain $\Omega_\vareps$ of the form \eqref{eq:Domain} consist of at most $N(d, k)$ branches of analytic functions which have at most algebraic singularities near $\vareps = 0$. At first-order in $\vareps$, the perturbation is given by the real eigenvalues of the Hermitian matrix $M^{(d, k)}$ of size $N(d, k)$, whose entries are given by \eqref{eq:Oeps_matrix}. 
\end{theorem} 

\begin{equation} \label{eq:Oeps_matrix1} 
M_{m, n}^{(d, k)} = \int_{S^d} \rho(\hat\theta)\left(-k(k + d)Y_k^m(\hat\theta)\cc{Y_k^n(\hat\theta)} + \nabla_{S^d}Y_k^m(\hat\theta)\cdot\nabla_{S^d}\cc{Y_k^n(\hat\theta)}\right) d\sigma_d. 
\end{equation} 

\begin{gather} \label{eq:Oeps_matrix1b} 
\begin{aligned} 
& -\sum_{j = 1}^d \int_{S^d} \frac{\del_j\rho}{\eta_j^2}\del_jY_k^m\cdot \cc{Y_k^n}\, d\sigma_d \\ 
& = \sum_{j = 1}^d \int_{S^d} \rho\frac{\del_jY_k^m}{\eta_j}\cdot \frac{\del_j\cc{Y_k^n}}{\eta_j}\, d\sigma_d + \sum_{j = 1}^d \int_{S^d} \rho \left(\frac{\del_j\left(\sin^{j - 1}(\theta_j)\del_jY_k^m\right)}{\eta_j^2\sin^{j - 1}(\theta_j)}\right) \cc{Y_k^n}\, d\sigma_d \\ 
& = \int_{S^d} \rho\nabla_{S^d}Y_k^m\cdot \nabla_{S^d}\cc{Y_k^n}\, d\sigma_d + \int_{S^d} \rho\left(\Delta_{S^d}Y_k^m\right)\cc{Y_k^n}\, d\sigma_d \\ 
& \stackrel{\eqref{eq:SH_eigs}}{=} \int_{S^d} \rho\nabla_{S^d}Y_k^m\cdot \nabla_{S^d}\cc{Y_k^n}\, d\sigma_d - \int_{S^d} k(k + d - 1)\rho Y_k^m\cc{Y_k^n}\, d\sigma_d. 
\end{aligned} 
\end{gather} 

\begin{lemma} \label{thm:Matrix_triple}
Given $d\ge 3$ and $k\in\Z^+$, let $M^{(d, k)}$ and $\rho$ be the Hermitian matrix and the perturbation function defined in Theorem \ref{thm:Oeps_matrix} and \eqref{eq:Domain}, respectively. Then the entries of $M^{(d, k)}$ can be written as 
\begin{equation} \label{eq:Matrix_triple} 
M_{m ,n}^{(d, k)} = -\frac{1}{2}\sum_{p = 0}^\infty\sum_{q = 1}^{N(d, k)} A_{p, q}\Big(p(p + d - 1) + 2k\Big)W_{q, m, n}^{p, k}, 
\end{equation} 
where 
\begin{equation*} 
W_{q, m, n}^{p, k} = \int_{S^d} Y_{p, q}(\hat\theta)Y_k^m(\hat\theta)\cc{Y_k^n(\hat\theta)}\, d\sigma_d. 
\end{equation*} 
\end{lemma} 

\begin{proof} 
From Theorem \ref{thm:Oeps_matrix} and \eqref{eq:SCgrad}, we have 
\begin{equation} \label{eq:Matrix_triple1} 
M_{m, n}^{(d, k)} \stackrel{\eqref{eq:Oeps_matrix1}}{=} -\int_{S^d} k\rho Y_k^m\cc{Y_k^n}\, d\sigma_d - \sum_{j = 1}^d\int_{S^d} \frac{\del_j\rho}{\eta_j^2}\del_jY_k^m\cdot \cc{Y_k^n}\, d\sigma_d. 
\end{equation}  
Integrating by parts with respect to $\theta_j$ and noting that $\eta_j$ is independent of $\theta_j$, we have that 
\begin{gather} 
\begin{aligned} \label{eq:Matrix_triple2}
-\sum_{j = 1}^d \int_{S^d} \frac{\del_j\rho}{\eta_j^2}\del_jY_k^m\cdot \cc{Y_k^n}\, d\sigma_d & = -\sum_{j = 1}^d \int_{S^{d - 1}}\int_{\theta_j} \frac{\del_jY_k^m}{\eta_j^2}\Big(\sin^{j- 1}(\theta_j)\del_j\rho\cdot \cc{Y_k^n}\Big)\, d\theta_j\, d\sigma_{d - 1, j} \\ 
& = \sum_{j = 1}^d \int_{S^{d - 1}}\int_{\theta_j} \frac{Y_k^m}{\eta_j^2}\del_j\Big(\sin^{j- 1}(\theta_j)\del_j\rho\cdot \cc{Y_k^n}\Big)\, d\theta_j\, d\sigma_{d - 1, j} \\ 
& = \sum_{j = 1}^d \int_{S^d} \frac{Y_k^m}{\eta_j^2}\left[\del_j\rho\cdot\del_j\cc{Y_k^n} + \frac{\del_j\left(\sin^{j - 1}(\theta_j)\del_j\rho\right)}{\sin^{j - 1}(\theta_j)}\cc{Y_k^n}\right] d\sigma_d \\ 
& \stackrel{\eqref{eq:SCbeltrami}}{=} \left(\sum_{j = 1}^d\int_{S^d} \frac{\del_j\rho}{\eta_j^2} Y_k^m\del_j\cc{Y_k^n}\, d\sigma_d\right) + \int_{S^d} \left(\Delta_{S^d}\rho\right) Y_k^m\cc{Y_k^n}\, d\sigma_d, 
\end{aligned} 
\end{gather} 
where the boundary term vanishes due to (1) $\del_1\rho$ and hyperspherical harmonics are $2\pi$-periodic with respect to $\theta_1$ for $j = 1$, and (2) $\sin^{j - 1}(0) = \sin^{j - 1}(\pi) = 0$ for $j = 2, 3, \dots, d$. On the other hand, we deduce from \eqref{eq:Oeps_matrix1b} from the proof of Theorem \ref{thm:Oeps_matrix} that 
\begin{equation} \label{eq:Matrix_triple3} 
-\sum_{j = 1}^d \int_{S^d} \frac{\del_j\rho}{\eta_j^2}\del_jY_k^m\cdot \cc{Y_k^n}\, d\sigma_d = -\sum_{j = 1}^d\int_{S^d} \frac{\del_j\rho}{\eta_j^2} Y_k^m\del_j\cc{Y_k^n}\, d\sigma_d. 
\end{equation} 
Taking the average of \eqref{eq:Matrix_triple2} and \eqref{eq:Matrix_triple3}, it follows that 
\begin{equation} \label{eq:Matrix_triple4}
-\sum_{j = 1}^d \int_{S^d} \frac{\del_j\rho}{\eta_j^2}\left(\del_jY_k^m\right)\cc{Y_k^n}\, d\sigma_d = \frac{1}{2}\int_{S^d} \left(\Delta_{S^d}\rho\right) Y_k^m\cc{Y_k^n}\, d\sigma_d. 
\end{equation} 
Finally, we substitute \eqref{eq:Matrix_triple4} and the expansion \eqref{eq:Domain} for $\rho$ into \eqref{eq:Matrix_triple1} to obtain 
\begin{align*} 
M_{m, n}^{(d, k)} & = -\frac{1}{2}\int_{S^d} \left(-\Delta_{S^d}\rho + 2k\rho\right)Y_k^m\cc{Y_k^n}\, d\sigma_d \\ 
& = -\frac{1}{2}\sum_{p = 0}^\infty \sum_{q = 0}^{N(d, p)} A_{p, q}\int_{S^d} \left(-\Delta_{S^d}Y_{p, q} + 2kY_{p, q}\right)Y_k^m\cc{Y_k^n}\, d\sigma_d \\ 
& \stackrel{\eqref{eq:SH_eigs}}{=} -\frac{1}{2}\sum_{p = 0}^\infty \sum_{q = 0}^{N(d, p)} \Big(p(p + d - 1) + 2k\Big)\int_{S^d} Y_{p, q}Y_k^m\cc{Y_k^n}\, d\sigma_d, 
% & = -\frac{1}{2}\sum_{p = 0}^\infty \sum_{q = 0}^{N(d, p)} \Big(p(p + d - 1) + 2k\Big)W_{q, m, n}^{p, k}, 
\end{align*} 
which gives the desired result. 
\end{proof} 

\end{document}
